\honest{If You Demand Magic, Magic Won't Help}{Eliezer Yudkowsky}{2008}

I open with Terry Pratchett: witches do not believe in the gods, because they know them too well. ``It would be like believing in the postman.''

I was raised on fantasy and I write it, so do not tell me I do not understand the genre. It occurred to me that if you could meet a fire-breathing dragon at the zoo while no one had ever seen a zebra, our fantasy stories would be full of zebras. ``The grass is always greener on the other side of unreality.''

In a standard plot, a good-hearted person from our world lands in a world of magic and becomes a ``(superpowerful)'' sorcerer. Readers presumably picture themselves in that role, and most of them are not scientists. ``Born into a world of science, they did not become scientists. What makes them think that, in a world of magic, they would act any differently?''

If they lack the capacity to be interested in merely real things, magic would not help them: once real, it would lose ``the charm of unattainability,'' like ``the lottery winners who, six months later, aren't nearly as happy as they expected to be.'' The excitement would wear off, probably as soon as they had to study spells. \nb{The best-known study of lottery winners, by Brickman, Coates and Janoff-Bulman (1978), found them no happier than controls and less pleased by mundane events. Its abstract gives no six-month figure and no measure of what the winners themselves had expected. A larger Swedish study (2020) found winners' life satisfaction raised for more than a decade, with smaller effects on happiness.} I do not say how many readers lack that capacity. Unless they can be as excited by hang-gliding as by riding a dragon, and by making light with electricity as by making it with magic, real dragons would excite them no more than hang-gliding does. I am not dissing dragons; we might even create some one day.

The same goes for the Future. Things do seem to get better, but this is the Future relative to the Dark Ages, with opportunities undreamt of by kings. If you reach the Future, you will find another Now. If your emotional energy can go only into a better tomorrow, no amount of time can help you. A future pill might fix even that, but my point is about which pills we should want to take.

A commenter, Matthew C., is excited by Rupert Sheldrake's theory, which ``explains'' protein folding and snowflake symmetry, phenomena I call ``non-explanation-demanding.'' Why is the commenter not just as excited by Special Relativity, which is known to be a law, so that the excitement is not wasted? ``If Sheldrake's theory were accepted truth taught in elementary schools, Matthew C. wouldn't care about it.'' \nb{The post gives no evidence for this beyond the commenter's interest in the theory. The comment itself gave a reason: that reductionism ``is incapable of explaining the real world.''} In the same way, the worst catastrophe for the New Age community would be for its rituals to start working and UFOs to appear: in a world where psychic powers were merely real, no one would believe in them, any more than anyone cares enough about gravity to believe in it. \nb{As with the commenter, the post states this as fact and gives no evidence about New Agers.}

I am not against magic. I am against a psychology that, in a world of spells, would pine for assembly lines. Binding yourself to reality means coming to terms with the fact that you live here; only then can you see your world's opportunities without wishing your sight away. I have found plenty of dragons to fight and magics to master in this world, and if I were transported into a fantasy novel I would expect to find myself studying the forbidden ultimate sorcery, because a new world would not change anything. ``It's not where you are, it's who you are.''

Last, the Litany Against Being Transported Into An Alternate Universe: if I am going to be happy, achieve greatness, learn true secrets, save the world, feel strongly or help people anywhere, ``I may as well do it in reality.''
